\begin{figure}[ht]
  \caption{Average Organizational Resilience}
    \label{fig:event_study}
    \centering
    \begin{minipage}[b]{0.325\textwidth}
        \centering
        \subcaption{Pull requests opened} \label{fig:event_study_prs_opened}
        \includegraphics[width=\textwidth]{\EventStudyDir/full_sample/norm/pull_request_opened.png}
    \end{minipage}
    \hfill
    \begin{minipage}[b]{0.325\textwidth}
        \centering
        \subcaption{Pull requests merged} \label{fig:event_study_prs_merged}
        \includegraphics[width=\textwidth]{\EventStudyDir/full_sample/norm/pull_request_merged.png}
    \end{minipage}
    \hfill
    \begin{minipage}[b]{0.325\textwidth}
        \centering
        \subcaption{New software releases}\label{fig:event_study_releases}
        \includegraphics[width=\textwidth]{\EventStudyDir/full_sample/norm/overall_new_release_count.png}
    \end{minipage}

  \begin{minipage}{1\textwidth}
    \textbf{Figure notes:} Panel~\subref{fig:event_study_prs_opened} presents event study estimates of the impact of key member departures on the number of pull requests opened.
    \EventStudySharedNotes
    The reported pre-trend p-value is from a Wald test of the null hypothesis that all pre-period coefficients equal zero. 
    Panels~\subref{fig:event_study_prs_merged} and~\subref{fig:event_study_releases} repeat this analysis using the number of pull requests merged and the number of new software releases as outcomes.

    
\end{minipage}
  

\end{figure}
